\begin{textsample}{POS Dim 1 – al – Score 48.00 – al/al\_em\_17.txt}  \label{ex:f1_pos_154}
ELEMENTOS DAS PRÁTICAS DIDÁTICO-PEDAGÓGICAS NA ÁREA Os componentes curriculares Arte, Educação Física, Língua \textbf{Inglesa} e Língua Portuguesa, que compõem a área de Linguagens e suas Tecnologias, são desenvolvidos considerando a construção do protagonismo, numa perspectiva inter e transdisciplinar de ensino e aprendizagem que prioriza o Projeto de Vida dos estudantes e a Educação Integral.
As práticas pedagógicas realizadas por esses componentes estão relacionadas por meio do que é comum à área : a linguagem.
Tais práticas \textbf{baseiam} -se nas metodologias ativas, entendidas como recursos a serviço dos docentes e através das quais são explorados : ensino híbrido ( sala de \textbf{aula} invertida, estação por rotação, laboratório rotacional, rotação individual ), estudo de caso, gamificação, aprendizagem \textbf{baseada} em \textbf{problemas}, aprendizagem \textbf{baseada} em projetos, atividades maker, aprendizado \textbf{baseado} em jogos, dentre outros.
As experiências fortalecem o trabalho coletivo e a autonomia das \textbf{juventudes} no Ensino Médio, sendo importante destacar que as ações pedagógicas são fundamentadas em sete \textbf{competências} específicas da área, que asseguram os direitos de aprendizagens dos estudantes e seguem \textbf{detalhadas} nos próximos \textbf{parágrafos}.
A \textbf{competência} 1 apresenta cinco habilidades a ela relacionadas e trata da combinação híbrida de diversas linguagens, a partir do trabalho com textos \textbf{complexos} e semióticos, por meio das possibilidades oferecidas pelas \textbf{tecnologias} \textbf{digitais}.
Nesse sentido, destaca que é necessário : 1.
Compreender o funcionamento das diferentes linguagens e práticas culturais ( artísticas, corporais e verbais ) e \textbf{mobilizar} esses conhecimentos na \textbf{recepção} e produção de discursos nos diferentes campos de atuação social e nas diversas mídias, para ampliar as formas de participação social, o entendimento e as possibilidades de explicação e interpretação \textbf{crítica} da realidade e para continuar aprendendo. ( BNCC, 2018, p. 483 ) Já o diálogo entre as práticas de linguagem artísticas, corporais e verbais, sem destaque aos saberes relacionados a um componente específico, evidencia -se na \textbf{competência} específica 2, que \textbf{salienta} a importância de : 2.
Compreender os processos identitários, conflitos e relações de poder que \textbf{permeiam} as práticas sociais de linguagem, respeitando as diversidades e a \textbf{pluralidade} de ideias e posições, e atuar socialmente com base em princípios e valores assentados na democracia, na igualdade e nos Direitos Humanos, exercitando o autoconhecimento, a empatia, o diálogo, a resolução de conflitos e a cooperação, e combatendo preconceitos de qualquer natureza. ( BNCC, 2018, p. 484 ) Essa \textbf{competência} aborda a \textbf{análise} da produção de sentidos envolvidos nos diferentes contextos do âmbito das práticas sociais da linguagem, por meio do desenvolvimento de quatro habilidades.
Esse conjunto de habilidades aponta para a necessidade de que o estudante do Ensino Médio deve compreender os variados processos de produção social da linguagem para interpretar as diversas produções artísticas, corporais e verbais existentes.
A \textbf{competência} específica 3, por sua vez, está relacionada a cinco habilidades, assim como a \textbf{competência} 1.
Tal \textbf{competência} dá ênfase à construção da autonomia do estudante, por meio da \textbf{análise} de processos de recuperação, \textbf{compreensão} e produção de discursos para uma constituição \textbf{crítica}, ética e relacionada aos Direitos Humanos, à consciência \textbf{socioambiental} e ao consumo responsável nos diferentes âmbitos. 3.
Utilizar diferentes linguagens ( artísticas, corporais e verbais ) para exercer, com autonomia e colaboração, protagonismo e \textbf{autoria} na vida pessoal e coletiva, de forma \textbf{crítica}, criativa, ética e solidária, defendendo pontos de \textbf{vista} que respeitem o outro e promovam os Direitos Humanos, a consciência \textbf{socioambiental} e o consumo responsável, em âmbito local, regional e global. ( BNCC, 2018, p. 485 ) As \textbf{competências} específicas 4, 5 e 6 abordam, com predominância, saberes historicamente construídos por componentes específicos.
Ao desenvolver a \textbf{competência} 4, o estudante deve : 4.
Compreender as línguas como fenômeno (geo)político, histórico, cultural, social, variável, heterogêneo e sensível aos contextos de uso, reconhecendo suas variedades e vivenciando -as como formas de expressões \textbf{identitárias}, pessoais e coletivas, bem como agindo no enfrentamento de preconceitos de qualquer natureza. ( BNCC, 2018, p. 486 ) Dessa forma, compreende -se o funcionamento das línguas e a adequação às situações de produção dos discursos.
Com essa finalidade, os estudantes, por meio de três habilidades, podem compreender a heterogeneidade e a variedade envolvidas nos processos de constituição dos sentidos, assim como a necessidade de seleção e organização de recursos estilísticos.
Nesse processo, destaca -se o respeito à variedade linguística sem preconceito.
Dessa forma, pode -se perceber que os conhecimentos relacionados aos componentes de Língua Portuguesa e Língua \textbf{Inglesa} destacam -se nessa \textbf{competência} por abordar predominantemente questões linguísticas.
Assim como a \textbf{competência} específica 4 aponta saberes relacionados às línguas, a \textbf{competência} 5 propõe a reflexão sobre preconceitos e estereótipos relacionados às práticas corporais, tais como os esportes, os jogos, as brincadeiras, as lutas e as artes marciais, as questões relacionadas à estética, à performance, à saúde, às manifestações culturais, ao lazer, ao entretenimento e, principalmente, às questões de gênero e sexualidade.
Dessa \textbf{maneira}, visa garantir as diversas expressões individuais corporais, \textbf{identitárias}, comportamentais e performáticas, a fim de prevenir e combater qualquer forma de preconceito e discriminação.
Para isso, o estudante deve compreender essas práticas e sua relação com fatores de ordem social, \textbf{econômica}, política e histórica, como também analisar os discursos que circulam sobre elas, por meio da leitura e da produção de discursos nas diferentes linguagens.
Percebe -se, então, que há ênfase às práticas corporais e aos saberes historicamente construídos pelo componente Educação Física, mas as outras práticas de linguagem, o que é comum à área de Linguagens e suas Tecnologias, também são contempladas na \textbf{análise} proposta pelas três habilidades dessa \textbf{competência}. 5.
Compreender os processos de produção e negociação de sentidos nas práticas corporais, reconhecendo -as e vivenciando -as como formas de expressão de valores e identidades, em uma perspectiva democrática e de respeito à diversidade. ( BNCC, 2018, p. 487 ) Por fim, a \textbf{competência} específica 6 prevê que o estudante, ao \textbf{final} do Ensino Médio, seja capaz de : 6.
Apreciar esteticamente as mais diversas produções artísticas e culturais, considerando suas características locais, regionais e globais, e \textbf{mobilizar} seus conhecimentos sobre as linguagens artísticas para dar significado e (re)construir produções autorais individuais e coletivas, exercendo protagonismo de \textbf{maneira} \textbf{crítica} e criativa, com respeito à diversidade de saberes, identidades e culturas. ( BNCC, 2018, p. 488 ) Nessa \textbf{competência}, há ênfase na \textbf{análise} e apreciação de critérios estéticos, relacionados aos diferentes contextos e movimentos culturais, históricos, entre outros.
Há, ainda, a previsão de criação de práticas artísticas, a partir dos repertórios pessoais em busca da ampliação do repertório do estudante.
Porém, nesse processo, também se pode perceber que há a intersecção com outras práticas de linguagens ( verbais e corporais ).
Finalmente, a \textbf{competência} específica 7, semelhante às \textbf{competências} 1, 2 e 3, prevê igualmente a \textbf{abordagem} de práticas artísticas, verbais e corporais.
Através do desenvolvimento de quatro habilidades, nessa \textbf{competência} deve -se : 7. \textbf{Mobilizar} práticas de linguagem no universo \textbf{digital}, considerando as dimensões técnicas, \textbf{críticas}, criativas, éticas e estéticas, para expandir as formas de produzir sentidos, de engajar -se em práticas autorais e coletivas, e de aprender a aprender nos campos da \textbf{ciência}, cultura, trabalho, informação e vida pessoal e coletiva. ( BNCC, 2018, p. 489 ) A partir das habilidades apresentadas pela \textbf{competência} 7, os estudantes serão capazes não apenas de utilizar as \textbf{tecnologias}, mas também devem desenvolver um \textbf{posicionamento} \textbf{crítico}, criativo, ético e estético em relação às práticas de linguagem do universo \textbf{digital}.
Na \textbf{competência} é prevista a exploração de recursos \textbf{digitais} para formação \textbf{crítica} e ética na produção e na circulação dessas linguagens.
Em síntese, na área de Linguagens e suas Tecnologias, três \textbf{competências} abordam com predominância conhecimentos historicamente relacionados a saberes específicos, divididos entre os seguintes componentes : Línguas ( \textbf{Competência} 4 ), Educação Física ( \textbf{Competência} 5 ) e Arte ( \textbf{Competência} 6 ).
As demais \textbf{competências} ( de número 1, 2, 3 e 7 ) atravessam de forma equilibrada todos os quatro componentes da área ( BRASIL, 2018, p. 489 ).
É importante \textbf{salientar}, ainda, que, mesmo havendo especificidades de \textbf{competências} em relação a determinado componente, é possível estabelecer integração entre os demais, porque a intersecção entre as diversas práticas de linguagem é fator intrínseco à proposição de todas as \textbf{competências}.
No \textbf{organizador} curricular da área, por exemplo, pode -se \textbf{verificar} essa intersecção.
Qualquer exemplo de prática pedagógica apresenta as três práticas de linguagem, mesmo que algumas habilidades indiquem a predominância já apontada anteriormente.
A seguir, apresenta -se um exemplo de prática pedagógica como sugestão para o desenvolvimento da segunda habilidade da \textbf{competência} 5 ( EM13LGG502 — Analisar criticamente preconceitos, estereótipos e relações de poder presentes nas práticas corporais, adotando \textbf{posicionamento} contrário a qualquer manifestação de injustiça e desrespeito a direitos humanos e valores democráticos ) : Fazer uma enquete na comunidade sobre músicas e/ou danças que destacam a \textbf{figura} da mulher e o impacto dessas expressões na cultura popular.
As músicas devem ser ouvidas, cantadas e dançadas.
Levar a letra das músicas mais citadas na enquete para sala de \textbf{aula} e propor \textbf{análise} sobre o contexto de produção, os padrões trazidos pelos elementos verbais, perfil do público ouvinte e mensagem transmitida.
Observar o movimento corporal e compreender o que a dança expressa, caracterizar o campo artístico-literário, a cultura de massa, fazendo relação entre o que foi observado e o preconceito e as relações de poder envolvidos.
Por fim, solicitar a elaboração de textos opinativos curtos que destaquem o respeito ao direito do outro e o combate ao preconceito de gênero, para ser \textbf{divulgado} em diferentes espaços da unidade escolar ( corredores, acesso, salas de \textbf{aula}, pátio etc. ), em data específica escolhida pela turma ; pode -se ainda reescrever as letras das músicas e criar novas coreografias.
Nesse exemplo, observa -se que o estereótipo da \textbf{figura} da mulher é analisado por meio de elementos verbais ( letras de músicas ), de movimentos corporais ( coreografias, dança ) e de manifestações artísticas ( música, dança ).
Dessa forma, é garantida a intersecção entre os diversos componentes da área por meio das práticas de linguagem.
É preciso destacar, ainda, que as sete \textbf{competências} específicas da área de Linguagens e suas Tecnologias relacionam -se a vinte e oito habilidades da área, no total.
Para o componente de Língua Portuguesa, além das habilidades da área, são apresentadas cinquenta e quatro habilidades específicas do componente.
Cada habilidade da área ou específica do componente de Língua Portuguesa está relacionada a objetos do conhecimento.
Os objetos do conhecimento, nessa perspectiva, são compreendidos como a seleção de determinado aspecto do saber a ser conhecido, aprofundado, desenvolvido ou ainda referenciado na rede de conhecimentos que compõem o repertório humano sociocultural.
Esses objetos são elementos essenciais no desenvolvimento das aprendizagens propostas e, segundo a BNCC ( 2018, p. 28 ), são “ entendidos como conteúdos, conceitos e processos ”.
É muito importante \textbf{salientar} que as próprias habilidades, em seu texto, apontam diretamente para esses objetos de conhecimento e que eles permitem a materialização das habilidades, como se pode \textbf{verificar} no exemplo a seguir : Nesse exemplo, há apresentação da \textbf{competência} específica de área de número 3, seguida da habilidade da área de número 04, relacionada a tal \textbf{competência}, e na \textbf{sequência} a explicitação de objetos de conhecimento já trazidos pela habilidade : Direitos Humanos, consciência \textbf{socioambiental} e consumo responsável.
Com o objetivo de desenvolver um trabalho com determinado objeto do conhecimento, o docente pode elencar os conteúdos relacionados a esse objeto e, de acordo com conceitos e propostas didático-pedagógicas, selecionar os processos de aprendizagem já acima referenciados. \textbf{COMPETÊNCIA} ESPECÍFICA DE ÁREA : 3.
Utilizar diferentes linguagens ( artísticas, corporais e verbais ) para exercer, com autonomia e colaboração, protagonismo e \textbf{autoria} na vida pessoal e coletiva, de forma \textbf{crítica}, criativa, ética e solidária, defendendo pontos de \textbf{vista} que respeitem o outro e promovam os Direitos Humanos, a consciência \textbf{socioambiental} e o consumo responsável, em âmbito local, regional e global.
HABILIDADE : ( EM13LGG304 ) Formular propostas, intervir e tomar decisões que levem em conta o bem comum e os Direitos Humanos, a consciência \textbf{socioambiental} e o consumo responsável em âmbito local, regional e global.
OBJETO DE CONHECIMENTO : Direitos Humanos, consciência \textbf{socioambiental} e consumo responsável.
Para \textbf{explicitar} a relação entre \textbf{competência}, habilidade e objeto do conhecimento na área de Linguagens e suas Tecnologias, foi elaborado um Organizador Curricular específico da área que apresenta os \textbf{Desdobramentos} Didático-Pedagógicos ( DesDP ).
Estes, por sua vez, \textbf{configuram} -se como sugestões de \textbf{sequências} didáticas, ou seja, como uma possibilidade de trabalho com as \textbf{competências} e as habilidades apresentadas pela BNCC.
Por meio da construção dos DesDP, portanto, pretende -se exemplificar práticas pedagógicas que considerem as \textbf{competências} e suas relações com as habilidades da área de Linguagens e suas Tecnologias.
Os DesDP aqui apresentados não pretendem excluir todas as outras possibilidades de trabalhar os objetos de conhecimento, mas pontuam \textbf{caminhos} possíveis.
Apresenta -se abaixo o \textbf{quadro} 01/28 para exemplificar a estruturação do Organizador Curricular e dos DesDP.
Com esse exemplo, é possível constatar que os \textbf{quadros} que compõem o \textbf{organizador} curricular apresentam, além da \textbf{competência} específica da área, da habilidade e do objeto de conhecimento, sugestões de práticas interdisciplinares que, por meio do trabalho com o objeto do conhecimento, garantem o desenvolvimento da habilidade e da \textbf{competência} em questão.
Pode -se observar, ainda, que o \textbf{quadro} está organizado em função das práticas de linguagens verbais, artísticas e corporais e que qualquer exemplo aponta as três práticas de linguagem, mesmo que aprofunde o conhecimento relacionado a uma dessas práticas.
Assim, o objeto de conhecimento “ Processos de produção e circulação de discursos ” é trabalhado, na primeira coluna, por meio de \textbf{vídeos} com \textbf{entrevistas} e textos jornalístico-midiáticos disponíveis na rede ; na segunda, o DesDP aborda textos com diferentes perspectivas disponíveis em jornais, revistas, blogs, redes sociais, charges e editoriais ; na terceira, a proposta é desenvolvida com a utilização de letras de reggae ; e na quarta, observa -se que o DesDP \textbf{sugere} um projeto de saúde para a terceira idade, por meio de \textbf{palestras}, vídeo-minuto, registro fotográfico.
Em todos os DesDP \textbf{verifica} -se que o objeto de conhecimento pode ser trabalhado de diferentes formas a partir da \textbf{análise} e da produção de diferentes práticas de linguagem.
Com tais \textbf{encaminhamentos}, os DesDP consideram o contexto na produção dos saberes e estimulam a relação entre as culturas local, regional e global, visando ao protagonismo do estudante e sua inserção \textbf{digital}.
Essa estruturação corrobora para o desenvolvimento integral do estudante, que é fundamentado na Lei nº 8.069/1990, do Estatuto da Criança e do Adolescente ( ECA ), e na formação do sujeito \textbf{crítico} e \textbf{reflexivo}.
A prática vivenciada nessa \textbf{abordagem} possibilita a ampliação do repertório do estudante, ao incorporar ao ensino e aprendizagem a diversidade de linguagens e suas esferas \textbf{discursivas}, a utilização das mídias e das \textbf{tecnologias} e a \textbf{contextualização} territorial e global.
Esses campos estão divididos em Campo da vida pessoal, Campo artístico-literário, Campo das práticas de estudo e pesquisa, Campo jornalístico-midiático e Campo de atuação na vida pública.

% matched lemmas: abordagem, análise, aula, autoria, basear, caminho, ciência, competência, complexo, compreensão, configurar, contextualização, crítico, desdobramento, detalhar, digital, discursivo, divulgar, econômico, encaminhamento, entrevista, explicitar, figura, final, identitária, inglês, juventude, maneira, mobilizar, organizador, palestra, parágrafo, permear, pluralidade, posicionamento, problema, quadro, recepção, reflexivo, salientar, sequência, socioambiental, sugerir, tecnologia, verificar, vista, vídeo
\end{textsample}
